
\chapter{Motivation}
\label{chap:motivation}


\textcolor{red}{a short focused section or chapter that states the specific modeling problem:}

%classical affine models are too rigid, neural approaches lack arbitrage-freeness, this model bridges the gap. This is where you connect the theory to the problem.}



%--------------------------------------------------------
%\subsection{Observed Swap Curve}
%--------------------------------------------------------

%In practice, the term structure is observed through a discrete set of market-quoted instruments. In many fixed-income markets, par swap rates across maturities provide a liquid representation of the yield curve.

%Let
%\[
%\mathcal{T}=\{\tau_1,\dots,\tau_m\}
%\]
%denote the set of quoted maturities. The observed swap curve at time $t$ is then given by the vector
%\[S_t=\bigl(S(t,\tau_1),\dots,S(t,\tau_m)\bigr)^\top \in \mathbb{R}^m.\]

%Since swap rates are determined by the discount curve through \eqref{eq:swap_rate}, modeling the swap curve amounts to modeling the underlying discount curve together with the transformation linking the two. However, because market quotes are observed only at finitely many maturities, the discrete swap curve does not by itself determine a unique continuous discount or forward-rate curve. Constructing a continuous arbitrage-consistent term structure therefore requires additional structure, either through interpolation or through a parametric or factor-based model.

%--------------------------------------------------------
%\subsection{Empirical Factor Structure of Yield Curves}
%--------------------------------------------------------

%Although the term structure is formally infinite-dimensional, empirical studies show that the cross-sectional variation of yield curves can often be approximated well by a small number of factors. These factors are commonly interpreted as level, slope, and curvature, although the exact number of factors retained depends on the modeling objective.

%This empirical observation motivates the use of low-dimensional factor models in term structure analysis. Classical specifications such as the Nelson-Siegel model capture this structure explicitly, while affine term structure models generate yield curves from a finite-dimensional vector of state variables. More generally, it motivates representing the term structure through a low-dimensional latent state variable, so that the infinite-dimensional curve is summarized by a finite set of factors evolving over time.
